% ===== PRÉAMBULE CANONIQUE — à coller VERBATIM en tête de CHAQUE .tex =====
\documentclass[11pt,a4paper]{article}
\usepackage[utf8]{inputenc}\usepackage[T1]{fontenc}\usepackage{lmodern}
\usepackage[french]{babel}
\usepackage{amsmath,amssymb,mathtools}\usepackage{icomma}
\usepackage[version=4]{mhchem}          % \ce{...} : équations & formules
\usepackage{siunitx}\sisetup{locale=FR} % \qty{}{}, \si{}, \ang{}
\usepackage{chemfig}                    % \chemfig{...} : structures développées
\usepackage{xcolor}\usepackage{colortbl}
\usepackage{tikz}\usepackage{pgfplots}\pgfplotsset{compat=1.18}
\usepackage[most]{tcolorbox}\usepackage{enumitem}\usepackage{array,booktabs}
\usepackage[a4paper,margin=2.2cm]{geometry}
\usetikzlibrary{arrows.meta,calc,angles,quotes,positioning,patterns,backgrounds,shapes.geometric}
\definecolor{primary}{RGB}{222,49,99}\definecolor{primarydark}{RGB}{150,22,55}
\definecolor{defcol}{RGB}{0,114,178}\definecolor{methcol}{RGB}{52,92,148}
\definecolor{excol}{RGB}{0,135,90}\definecolor{attncol}{RGB}{200,40,55}
\tcbset{boxbase/.style={enhanced,breakable,boxrule=0.9pt,arc=2.5pt,
  left=3mm,right=3mm,top=2.3mm,bottom=2.3mm,fonttitle=\bfseries,coltitle=white}}
% --- boîtes TUTORIEL (noms EXACTS, ne pas renommer) ---
\newtcolorbox{cle}[1][]{boxbase,colback=defcol!6,colframe=defcol,
  title={\textsc{À retenir}\ifx\relax#1\relax\relax\else\textmd{ : \itshape #1}\fi}}
\newtcolorbox{methode}[1][]{boxbase,colback=methcol!7,colframe=methcol,sharp corners,
  title={\textsc{Méthode}\ifx\relax#1\relax\relax\else\textmd{ --- \itshape #1}\fi}}
\newtcolorbox{exemplebox}[1][]{boxbase,colback=excol!5,colframe=excol,
  title={\textsc{Exemple}\ifx\relax#1\relax\relax\else\textmd{ : \itshape #1}\fi}}
\newtcolorbox{piege}[1][]{boxbase,colback=attncol!6,colframe=attncol,
  title={\textsc{Piège}\ifx\relax#1\relax\relax\else\textmd{ : \itshape #1}\fi}}
\newtcolorbox{remarque}[1][]{boxbase,colback=black!4,colframe=black!45,
  title={\textsc{Remarque}\ifx\relax#1\relax\relax\else\textmd{ : \itshape #1}\fi}}
% --- boîtes EXERCICE / CORRIGÉ (noms EXACTS, auto-comptées) ---
\newtcolorbox[auto counter]{exo}[1][]{boxbase,colback=primary!4,colframe=primary,
  title={\textsc{Exercice}~\thetcbcounter\ifx\relax#1\relax\relax\else\textmd{ --- \itshape #1}\fi}}
\newtcolorbox[auto counter]{corrige}[1][]{boxbase,colback=excol!4,colframe=excol,
  title={\textsc{Corrigé}~\thetcbcounter\ifx\relax#1\relax\relax\else\textmd{ --- \itshape #1}\fi}}
\newcommand{\R}{\mathbb{R}}\newcommand{\N}{\mathbb{N}}\newcommand{\Z}{\mathbb{Z}}\newcommand{\ds}{\displaystyle}
% (un \usepackage{fancyhdr}+en-tête est possible ; il est ignoré côté web)
% =========================================================================
\begin{document}

\begin{corrige}[composition d'un cation]
Protons $=Z$, neutrons $=A-Z$, électrons $=Z-q$.
\begin{enumerate}[label=\alph*)]
  \item $\ce{^{24}_{12}Mg^{2+}}$ : $n(p^+)=12$, $n(n^0)=24-12=12$, $n(e^-)=12-2=10$.
  \item $\ce{^{23}_{11}Na^{+}}$ : $n(p^+)=11$, $n(n^0)=23-11=12$, $n(e^-)=11-1=10$.
  \item $\ce{^{27}_{13}Al^{3+}}$ : $n(p^+)=13$, $n(n^0)=27-13=14$, $n(e^-)=13-3=10$.
\end{enumerate}
Les trois cations ont $10$ électrons (configuration du néon), mais des noyaux différents.
\end{corrige}

\begin{corrige}[composition d'un anion]
Protons $=Z$, neutrons $=A-Z$, électrons $=Z-q$ (avec $q<0$, on ajoute $|q|$ électrons).
\begin{enumerate}[label=\alph*)]
  \item $\ce{^{35}_{17}Cl^{-}}$ : $n(p^+)=17$, $n(n^0)=35-17=18$, $n(e^-)=17-(-1)=18$.
  \item $\ce{^{16}_{8}O^{2-}}$ : $n(p^+)=8$, $n(n^0)=16-8=8$, $n(e^-)=8-(-2)=10$.
  \item $\ce{^{32}_{16}S^{2-}}$ : $n(p^+)=16$, $n(n^0)=32-16=16$, $n(e^-)=16-(-2)=18$.
\end{enumerate}
\end{corrige}

\begin{corrige}[le noyau ne bouge pas]
\begin{enumerate}[label=\alph*)]
  \item \textbf{$0$ électron.} Le noyau ne contient que des nucléons (protons et neutrons) : aucun
        électron ne s'y trouve jamais. Les $18$ électrons de $\ce{Sc^{3+}}$ sont dans le nuage
        électronique.
  \item Seul $n(e^-)$ change ($20 \to 18$). Le noyau étant intact, $n(p^+)=20$ et $n(n^0)=20$ restent
        \emph{identiques}.
  \item Oui : même noyau $\Rightarrow$ même nombre de nucléons, $A=40$ pour les deux.
\end{enumerate}
\end{corrige}

\begin{corrige}[retrouver la charge d'un ion]
On calcule $q = Z - n(e^-)$ ($q>0$ : cation, $q<0$ : anion).
\begin{enumerate}[label=\alph*)]
  \item $q = 12 - 10 = +2$ : cation $\ce{^{24}_{12}Mg^{2+}}$.
  \item $q = 8 - 10 = -2$ : anion $\ce{^{16}_{8}O^{2-}}$.
  \item $q = 26 - 24 = +2$ : cation $\ce{^{56}_{26}Fe^{2+}}$.
\end{enumerate}
\end{corrige}

\begin{corrige}[atome contre ion : compter et comparer]
Le noyau est le même pour l'atome et son ion ; seul le nombre d'électrons diffère.
\begin{enumerate}[label=\alph*)]
  \item $\ce{^{23}_{11}Na}$ : $(11,\,12,\,11)$ ; \quad $\ce{^{23}_{11}Na^{+}}$ : $(11,\,12,\,10)$.
        Le cation a \emph{un électron de moins} ; protons et neutrons sont identiques.
  \item $\ce{^{32}_{16}S}$ : $(16,\,16,\,16)$ ; \quad $\ce{^{32}_{16}S^{2-}}$ : $(16,\,16,\,18)$.
        L'anion a \emph{deux électrons de plus} ; protons et neutrons sont identiques.
\end{enumerate}
\end{corrige}

\end{document}
